
\documentclass[12pt,a4paper]{article}

% ============================================================
% PACOTES
% ============================================================

\usepackage[utf8]{inputenc}
\usepackage[T1]{fontenc}
\usepackage[brazil]{babel}

\usepackage{amsmath,amssymb,amsfonts,mathtools}
\usepackage{bm}
\usepackage{booktabs}
\usepackage{geometry}
\usepackage{graphicx}
\usepackage{float}
\usepackage{hyperref}
\newcommand{\pd}[2]{\frac{\partial #1}{\partial #2}}
\usepackage{tikz}
\usetikzlibrary{calc, decorations.pathmorphing, arrows.meta}
\usepackage{amsmath}

\geometry{a4paper, margin=2.5cm}

\hypersetup{
    colorlinks=true,
    linkcolor=blue,
    citecolor=blue,
    urlcolor=blue
}

% ============================================================
% Estilos
% ============================================================

\tikzset{
    originalsite/.style={
        circle,
        fill=gray!55,
        inner sep=0pt,
        minimum size=2.5pt
    },
    blockbond/.style={
        draw=blue!75!black,
        line width=1.4pt
    },
    blockcenter/.style={
        circle,
        fill=red!75!black,
        draw=black,
        line width=0.4pt,
        inner sep=0pt,
        minimum size=5.5pt
    },
    stagearrow/.style={
        -{Latex[length=3mm]},
        line width=0.9pt,
        draw=black!70
    }
}


\begin{document}

% ============================================================
% CABEÇALHO
% ============================================================

\begin{center}
    {\Large \textbf{Quinta série de exercícios}}\\[0.6cm]
    {\large Tópicos de Mecânica Estatística -- Transições de Fases -- 2026}\\[0.2cm]
    {\large Grupo de Renormalização -- 05/10/2026}\\[0.2cm]
    {\large Data de Entrega -- 04/10/2026}\\[0.2cm]
    {\Large \textbf{Autor:} Thiago Siqueira Domingues (Nº USP: 8942194)}
\end{center}
%
\vspace{0.4cm}
% ============================================================
% PROBLEMA 1
% ============================================================
\section*{Problema 1 -- Modelo de Ising na rede hierárquica do diamante}
%
Consideramos o modelo de Ising ferromagnético em campo nulo,
%
\begin{equation}
    \mathcal H
    =
    -J\sum_{\langle ij\rangle}\sigma_i\sigma_j,
    \qquad
    J>0,
    \qquad
    \sigma_i=\pm1.
    \label{eq:H_ising_p1}
\end{equation}
%
Introduzimos o acoplamento adimensional
%
\begin{equation}
    K\equiv\beta J=\frac{J}{k_BT}.
    \label{eq:K_p1}
\end{equation}
%
A construção da rede do diamante substitui uma ligação por duas cadeias em paralelo, cada uma contendo duas ligações em série. Assim, o fator de escala linear é
%
\begin{equation}
    b=2,
\end{equation}
%
e o número efetivo de ramos em paralelo é também \(p=2\). A dimensão da
rede hierárquica é, portanto,
%
\begin{equation}
    d
    =
    \frac{\ln(pb)}{\ln b}
    =
    \frac{\ln 4}{\ln 2}
    =
    2.
    \label{eq:dim_diamond}
\end{equation}
% ------------------------------------------------------------
Agora, considere uma cadeia formada pelos spins extremos \(\sigma_1\) e
\(\sigma_2\) e pelo spin intermediário \(s\):
%
\begin{equation}
    \sigma_1-s-\sigma_2.
\end{equation}
%
O peso de Boltzmann dessa cadeia é
%
\begin{equation}
    W(\sigma_1,\sigma_2)
    =
    \sum_{s=\pm1}
    \exp\left[
        K s(\sigma_1+\sigma_2)
    \right].
    \label{eq:W_chain_p1}
\end{equation}
%
Analisamos separadamente os dois casos possíveis para os spins externos. Se \(\sigma_1=\sigma_2\), então \(\sigma_1+\sigma_2=2\sigma_1\). Como
\(\sigma_1=\pm1\), podemos escrever \(\sigma_1^2=1\). Assim,
%
\begin{equation}
    W_{\parallel}
    =
    \sum_{s=\pm1}
    e^{2K s\sigma_1}
    =
    e^{2K\sigma_1}+e^{-2K\sigma_1}.
\end{equation}
%
Como a soma sobre \(s=\pm1\) é simétrica, o resultado não depende do
sinal de \(\sigma_1\). Portanto,
%
\begin{equation}
    W_{\parallel}
    =
    e^{2K}+e^{-2K}
    =
    2\cosh(2K).
    \label{eq:Wpar_p1}
\end{equation}

\medskip
\noindent\textbf{Caso 2: \(\sigma_1=-\sigma_2\).}
%
Se \(\sigma_1=-\sigma_2\), então \(\sigma_1+\sigma_2=0\). Logo,
%
\begin{equation}
    W_{\mathrm{anti}}
    =
    \sum_{s=\pm1}
    e^{0}
    =
    \sum_{s=\pm1}1
    =
    2.
    \label{eq:Wanti_p1}
\end{equation}

\medskip
Queremos substituir a cadeia por uma ligação efetiva:
%
\begin{equation}
    W(\sigma_1,\sigma_2)
    =
    A\,e^{K_s\sigma_1\sigma_2}.
    \label{eq:W_eff_p1}
\end{equation}
%
Para \(\sigma_1=\sigma_2\), temos \(\sigma_1\sigma_2=1\), de modo que
%
\begin{equation}
    W_{\parallel}=A\,e^{K_s}.
\end{equation}
%
Para \(\sigma_1=-\sigma_2\), temos \(\sigma_1\sigma_2=-1\), de modo que
%
\begin{equation}
    W_{\mathrm{anti}}=A\,e^{-K_s}.
\end{equation}
%
Dividindo as duas equações,
%
\begin{equation}
    \frac{W_{\parallel}}{W_{\mathrm{anti}}}
    =
    \frac{A\,e^{K_s}}{A\,e^{-K_s}}
    =
    e^{2K_s}.
\end{equation}
%
Substituindo \eqref{eq:Wpar_p1} e \eqref{eq:Wanti_p1},
%
\begin{equation}
    e^{2K_s}
    =
    \frac{2\cosh(2K)}{2}
    =
    \cosh(2K).
    \label{eq:e2Ks_p1}
\end{equation}
%
Tomando o logaritmo,
%
\begin{equation}
    2K_s=\ln\cosh(2K),
\end{equation}
%
e, portanto,
%
\begin{equation}
    \boxed{
    K_s
    =
    \frac12\ln\cosh(2K).
    }
    \label{eq:Ks_p1}
\end{equation}

\medskip
Usando a identidade
%
\begin{equation}
    \operatorname{arctanh}(\tanh^2K)
    =
    \frac12\ln\cosh(2K),
    \label{eq:identity_arctanh_p1}
\end{equation}
%
também podemos escrever
%
\begin{equation}
    \boxed{
    K_s=\operatorname{arctanh}(\tanh^2K).
    }
    \label{eq:Ks_arctanh_p1}
\end{equation}
%
Para verificar \eqref{eq:identity_arctanh_p1}, seja
%
\begin{equation}
    t=\tanh K.
\end{equation}
%
Então
%
\begin{equation}
    \tanh^2K=t^2,
\end{equation}
%
e, pela definição de \(\operatorname{arctanh}\),
%
\begin{equation}
    \operatorname{arctanh}(t^2)
    =
    \frac12
    \ln\left(
        \frac{1+t^2}{1-t^2}
    \right).
    \label{eq:arctanh_def_p1}
\end{equation}
%
Por outro lado,
%
\begin{equation}
    \cosh(2K)
    =
    \frac{e^{2K}+e^{-2K}}{2}.
\end{equation}
%
Usando \(t=\tanh K\), temos
%
\begin{equation}
    e^{2K}=\frac{1+t}{1-t}.
\end{equation}
%
Portanto,
%
\begin{align}
    \cosh(2K)
    &=
    \frac12
    \left[
        \frac{1+t}{1-t}
        +
        \frac{1-t}{1+t}
    \right]
    \nonumber\\[4pt]
    &=
    \frac12
    \frac{
        (1+t)^2+(1-t)^2
    }{(1-t)(1+t)}
    \nonumber\\[4pt]
    &=
    \frac12
    \frac{
        1+2t+t^2+1-2t+t^2
    }{1-t^2}
    \nonumber\\[4pt]
    &=
    \frac12
    \frac{
        2+2t^2
    }{1-t^2}
    =
    \frac{1+t^2}{1-t^2}.
    \label{eq:cosh2K_t_p1}
\end{align}
%
Substituindo \eqref{eq:cosh2K_t_p1} em \eqref{eq:arctanh_def_p1},
%
\begin{equation}
    \operatorname{arctanh}(t^2)
    =
    \frac12
    \ln\left(
        \cosh(2K)
    \right)
    =
    \frac12\ln\cosh(2K),
\end{equation}
%
o que confirma \eqref{eq:identity_arctanh_p1}.
% ------------------------------------------------------------
Na rede do diamante existem dois ramos idênticos em paralelo. Os
acoplamentos efetivos de ramos em paralelo se somam, de modo que
%
\begin{equation}
    K'=2K_s.
    \label{eq:Kprime_parallel_p1}
\end{equation}
%
Substituindo \eqref{eq:Ks_p1},
%
\begin{equation}
    \boxed{
    K'
    =
    \ln\cosh(2K).
    }
    \label{eq:RG_Ising_diamond}
\end{equation}
%
Assim, a relação de recorrência é
%
\begin{equation}
    \boxed{
    K_{n+1}
    =
    \ln\cosh(2K_n).
    }
    \label{eq:recursion_p1}
\end{equation}
% 
Um ponto fixo satisfaz
%
\begin{equation}
    K^*=\ln\cosh(2K^*).
    \label{eq:fixed_p1}
\end{equation}
%
Existe imediatamente o ponto
%
\begin{equation}
    \boxed{K^*_0=0},
\end{equation}
%
que corresponde ao limite de alta temperatura, \(T\to\infty\), pois
\(K=\beta J\to0\). Existe também um ponto fixo não trivial. Para encontrá-lo,
escrevemos \eqref{eq:fixed_p1} na forma
%
\begin{equation}
    e^{K^*}=\cosh(2K^*).
    \label{eq:fixed_exp_p1}
\end{equation}
%
Resolvendo numericamente \eqref{eq:fixed_exp_p1}, obtemos
%
\begin{equation}
    \boxed{
    K_c\simeq0.6093779.
    }
    \label{eq:Kc_p1}
\end{equation}
%
Como
%
\begin{equation}
    K_c=\frac{J}{k_BT_c},
\end{equation}
%
a temperatura de transição estimada é
%
\begin{equation}
    \boxed{
    \frac{k_BT_c^{\mathrm{MK}}}{J}
    =
    \frac{1}{K_c}
    \simeq1.641.
    }
    \label{eq:Tc_p1}
\end{equation}
% 
Seja
%
\begin{equation}
    K'=R(K),
    \qquad
    R(K)=\ln\cosh(2K).
    \label{eq:RK_p1}
\end{equation}
%
Próximo ao ponto fixo,
%
\begin{equation}
    K=K_c+\delta K,
\end{equation}
%
e a transformação linearizada é
%
\begin{equation}
    \delta K'
    =
    R'(K_c)\,\delta K.
    \label{eq:linearized_p1}
\end{equation}
%
Calculamos \(R'(K)\):
%
\begin{align}
    R'(K)
    &=
    \frac{d}{dK}
    \ln\cosh(2K)
    \nonumber\\[4pt]
    &=
    \frac{1}{\cosh(2K)}
    \frac{d}{dK}
    \cosh(2K)
    \nonumber\\[4pt]
    &=
    \frac{1}{\cosh(2K)}
    \,
    2\sinh(2K)
    \nonumber\\[4pt]
    &=
    2\tanh(2K).
    \label{eq:Rprime_p1}
\end{align}
%
Portanto,
%
\begin{equation}
    \lambda_t
    =
    R'(K_c)
    =
    2\tanh(2K_c).
    \label{eq:lambdat_form_p1}
\end{equation}
%
Para calcular \(\tanh(2K_c)\), usamos a equação do ponto fixo
\eqref{eq:fixed_exp_p1}:
%
\begin{equation}
    \cosh(2K_c)=e^{K_c}.
\end{equation}
%
Então
%
\begin{equation}
    \sinh^2(2K_c)
    =
    \cosh^2(2K_c)-1
    =
    e^{2K_c}-1,
\end{equation}
%
e, como \(\sinh(2K_c)>0\),
%
\begin{equation}
    \sinh(2K_c)
    =
    \sqrt{e^{2K_c}-1}.
\end{equation}
%
Logo,
%
\begin{equation}
    \tanh(2K_c)
    =
    \frac{\sinh(2K_c)}{\cosh(2K_c)}
    =
    \frac{\sqrt{e^{2K_c}-1}}{e^{K_c}}
    =
    \sqrt{1-e^{-2K_c}}.
    \label{eq:tanh2Kc_p1}
\end{equation}
%
Substituindo \eqref{eq:tanh2Kc_p1} em \eqref{eq:lambdat_form_p1},
%
\begin{equation}
    \lambda_t
    =
    2\sqrt{1-e^{-2K_c}}.
    \label{eq:lambdat_sqrt_p1}
\end{equation}
%
Com \(K_c\simeq0.6093779\),
%
\begin{equation}
    e^{-2K_c}
    \simeq
    e^{-1.2187558}
    \simeq
    0.2956,
\end{equation}
%
de modo que
%
\begin{equation}
    \lambda_t
    \simeq
    2\sqrt{1-0.2956}
    =
    2\sqrt{0.7044}
    \simeq
    1.67857.
    \label{eq:lambdat_p1}
\end{equation}
%
Pela definição do expoente térmico,
%
\begin{equation}
    \lambda_t=b^{y_t}.
\end{equation}
%
Como \(b=2\),
%
\begin{equation}
    y_t
    =
    \frac{\ln\lambda_t}{\ln 2}
    \simeq
    \frac{\ln 1.67857}{\ln 2}
    \simeq
    0.74724.
    \label{eq:yt_p1}
\end{equation}
%
Assim,
%
\begin{equation}
    \boxed{
    y_t\simeq0.74724.
    }
    \label{eq:yt_box_p1}
\end{equation}
%
O expoente \(\nu\) é
%
\begin{equation}
    \nu=\frac{1}{y_t},
\end{equation}
%
dando
%
\begin{equation}
    \boxed{
    \nu_{\mathrm{MK}}\simeq1.33827.
    }
    \label{eq:nu_p1}
\end{equation}
%
Usando a relação de hiperescala em \(d=2\),
%
\begin{equation}
    2-\alpha=d\nu,
\end{equation}
%
e como \(d=2\),
%
\begin{equation}
    2-\alpha=2\nu.
\end{equation}
%
Portanto,
%
\begin{equation}
    \boxed{
    \alpha_{\mathrm{MK}}
    =
    2-2\nu
    \simeq-0.67653.
    }
    \label{eq:alpha_p1}
\end{equation}
% 
Para o modelo de Ising na rede quadrada, a solução exata de Onsager fornece
%
\begin{equation}
    K_c^{\mathrm{exato}}
    =
    \frac12\ln(1+\sqrt2)
    \simeq0.440687,
\end{equation}
%
ou
%
\begin{equation}
    \boxed{
    \frac{k_BT_c^{\mathrm{exato}}}{J}
    \simeq2.269185.
    }
\end{equation}
%
Além disso,
%
\begin{equation}
    \nu_{\mathrm{exato}}=1,
    \qquad
    \alpha_{\mathrm{exato}}=0.
\end{equation}
%
A comparação é apresentada na Tabela~\ref{tab:comparison_p1}.
%
\begin{table}[H]
    \centering
    \caption{Resultados para o modelo de Ising na rede quadrada.}
    \label{tab:comparison_p1}
    \begin{tabular}{lcc}
        \toprule
        Método
        & \(k_BT_c/J\)
        & \(\nu\) \\
        \midrule
        Migdal--Kadanoff
        & \(1.641\)
        & \(1.338\) \\
        Exato
        & \(2.269\)
        & \(1\) \\
        \bottomrule
    \end{tabular}
\end{table}
%
A transformação de Migdal--Kadanoff reproduz qualitativamente a estrutura de uma transição crítica, mas fornece valores aproximados para a temperatura crítica e para os expoentes. Para uma rede hierárquica construída com \(b\) bonds em série e \(p\) ramos em paralelo, o resultado geral para o modelo de Ising é
%
\begin{equation}
    K_s
    =
    \operatorname{arctanh}
    \left[
        (\tanh K)^b
    \right],
    \label{eq:Ks_general_p1}
\end{equation}
%
seguido por
%
\begin{equation}
    \boxed{
    K'
    =
    p\,\operatorname{arctanh}
    \left[
        (\tanh K)^b
    \right].
    }
    \label{eq:general_hierarchical}
\end{equation}
%
A dimensão correspondente é
%
\begin{equation}
    d
    =
    \frac{\ln(pb)}{\ln b}.
    \label{eq:dim_general_p1}
\end{equation}
%
No caso da rede do diamante, temos \(b=2\) e \(p=2\). Substituindo em
\eqref{eq:general_hierarchical},
%
\begin{equation}
    K'
    =
    2\,\operatorname{arctanh}
    \left[
        (\tanh K)^2
    \right].
\end{equation}
%
Usando a identidade \eqref{eq:identity_arctanh_p1},
%
\begin{equation}
    K'
    =
    2\cdot
    \frac12\ln\cosh(2K)
    =
    \ln\cosh(2K),
\end{equation}
%
que coincide com \eqref{eq:RG_Ising_diamond}. Isso mostra a consistência
da fórmula geral. Assim, para a ``necklace'' utilizada em aula, basta identificar os
valores de \(p\) e \(b\) de sua construção para obter a recorrência e,
consequentemente, seu ponto fixo e expoente térmico.
% ============================================================
% PROBLEMA 2
% ============================================================
\section*{Problema 2 -- Modelo de Potts na rede hierárquica do diamante}
%
Considere o modelo de Potts com \(q\ge2\) estados:
%
\begin{equation}
    \mathcal H
    =
    -J\sum_{\langle ij\rangle}
    \delta_{\sigma_i,\sigma_j},
    \qquad
    \sigma_i=0,1,\ldots,q-1,
    \qquad
    J>0.
    \label{eq:H_potts}
\end{equation}
%
A construção da rede do diamante é a mesma do Problema~1: uma ligação
é substituída por dois ramos em paralelo, cada um contendo dois bonds
em série. Assim, temos \(p=2\) ramos em paralelo e \(b=2\) bonds em
série. A dimensão da rede hierárquica é
%
\begin{equation}
    d=\frac{\ln(pb)}{\ln b}
    =\frac{\ln 4}{\ln 2}=2.
\end{equation}
% 
Considere os estados externos \(\sigma_1\) e \(\sigma_2\) e um estado intermediário \(\tau\). O peso de Boltzmann do ramo de dois bonds é
%
\begin{equation}
    W(\sigma_1,\sigma_2)
    =
    \sum_{\tau=0}^{q-1}
    e^{K\delta_{\sigma_1,\tau}}
    e^{K\delta_{\tau,\sigma_2}}.
    \label{eq:Wpotts}
\end{equation}
%
Analisamos separadamente os casos \(\sigma_1=\sigma_2\) e
\(\sigma_1\neq\sigma_2\). Se \(\tau=\sigma_1=\sigma_2\), então
%
\begin{equation}
    \delta_{\sigma_1,\tau}=1,
    \qquad
    \delta_{\tau,\sigma_2}=1,
\end{equation}
%
e o peso desse termo é
%
\begin{equation}
    e^{K\cdot1}e^{K\cdot1}=e^{2K}.
\end{equation}
%
Para os outros \(q-1\) estados de \(\tau\), temos
%
\begin{equation}
    \tau\neq\sigma_1,
    \qquad
    \tau\neq\sigma_2,
\end{equation}
%
de modo que
%
\begin{equation}
    \delta_{\sigma_1,\tau}=0,
    \qquad
    \delta_{\tau,\sigma_2}=0,
\end{equation}
%
e cada um desses termos contribui com peso
%
\begin{equation}
    e^{0}e^{0}=1.
\end{equation}
%
Portanto,
%
\begin{equation}
    W_{\parallel}
    =
    e^{2K}+(q-1).
    \label{eq:Wpar_potts}
\end{equation}
%
Se \(\tau=\sigma_1\), então
%
\begin{equation}
    \delta_{\sigma_1,\tau}=1,
    \qquad
    \delta_{\tau,\sigma_2}=0,
\end{equation}
%
pois \(\sigma_1\neq\sigma_2\). O peso é
%
\begin{equation}
    e^{K}e^{0}=e^{K}.
\end{equation}
%
Se \(\tau=\sigma_2\), então
%
\begin{equation}
    \delta_{\sigma_1,\tau}=0,
    \qquad
    \delta_{\tau,\sigma_2}=1,
\end{equation}
%
e o peso é novamente
%
\begin{equation}
    e^{0}e^{K}=e^{K}.
\end{equation}
%
Para os \(q-2\) estados restantes, temos
%
\begin{equation}
    \tau\neq\sigma_1,
    \qquad
    \tau\neq\sigma_2,
\end{equation}
%
e cada um contribui com peso \(1\). Logo,
%
\begin{equation}
    W_{\mathrm{anti}}
    =
    2e^{K}+(q-2).
    \label{eq:Wanti_potts}
\end{equation}
%
Representamos o ramo efetivo por uma ligação de Potts:
%
\begin{equation}
    W_{\mathrm{eff}}
    =
    A\,e^{K_s\delta_{\sigma_1,\sigma_2}}.
\end{equation}
%
Assim,
%
\begin{equation}
    W_{\parallel}=A\,e^{K_s},
    \qquad
    W_{\mathrm{anti}}=A.
\end{equation}
%
Dividindo as duas expressões,
%
\begin{equation}
    e^{K_s}
    =
    \frac{W_{\parallel}}{W_{\mathrm{anti}}}
    =
    \frac{e^{2K}+q-1}{2e^{K}+q-2}.
\end{equation}
%
Tomando o logaritmo,
%
\begin{equation}
    K_s
    =
    \ln
    \left[
        \frac{e^{2K}+q-1}{2e^{K}+q-2}
    \right].
    \label{eq:Ks_potts}
\end{equation}
% 
Na rede do diamante existem dois ramos idênticos em paralelo. Os
acoplamentos efetivos de ramos em paralelo se somam, de modo que
%
\begin{equation}
    K'=2K_s.
\end{equation}
%
Substituindo \eqref{eq:Ks_potts},
%
\begin{equation}
    \boxed{
    K'
    =
    2\ln
    \left[
        \frac{e^{2K}+q-1}{2e^{K}+q-2}
    \right].
    }
    \label{eq:potts_RG}
\end{equation}
%
É conveniente introduzir a variável
%
\begin{equation}
    v=e^{K}-1,
    \label{eq:v_def_potts}
\end{equation}
%
de modo que
%
\begin{equation}
    e^{K}=1+v.
\end{equation}
%
Calculamos separadamente os termos que aparecem em \eqref{eq:potts_RG}:
%
\begin{align}
    e^{2K}
    &=
    (e^{K})^2
    =
    (1+v)^2
    =
    1+2v+v^2,
    \label{eq:e2K_potts}
    \\[4pt]
    e^{2K}+q-1
    &=
    (1+2v+v^2)+q-1
    =
    v^2+2v+q,
    \label{eq:num_potts}
    \\[4pt]
    2e^{K}+q-2
    &=
    2(1+v)+q-2
    =
    q+2v.
    \label{eq:den_potts}
\end{align}
%
Portanto,
%
\begin{equation}
    K'
    =
    2\ln
    \left(
        \frac{v^2+2v+q}{q+2v}
    \right)
    =
    \ln
    \left(
        \frac{v^2+2v+q}{q+2v}
    \right)^2.
    \label{eq:Kprime_v_potts}
\end{equation}
%
Como
%
\begin{equation}
    e^{K'}=1+v',
\end{equation}
%
segue de \eqref{eq:Kprime_v_potts} que
%
\begin{equation}
    1+v'
    =
    \left(
        \frac{v^2+2v+q}{q+2v}
    \right)^2.
    \label{eq:1vprime_potts}
\end{equation}
%
Subtraindo \(1\) de ambos os lados,
%
\begin{equation}
    v'
    =
    \left(
        \frac{v^2+2v+q}{q+2v}
    \right)^2
    -
    1
    =
    \frac{
        (v^2+2v+q)^2-(q+2v)^2
    }{(q+2v)^2}.
    \label{eq:vprime_unsimplified_potts}
\end{equation}
%
Fatoramos o numerador usando a diferença de quadrados:
%
\begin{equation}
    (v^2+2v+q)^2-(q+2v)^2
    =
    \left[
        (v^2+2v+q)-(q+2v)
    \right]
    \left[
        (v^2+2v+q)+(q+2v)
    \right].
    \label{eq:diff_squares_potts}
\end{equation}
%
Calculamos cada fator:
%
\begin{align}
    (v^2+2v+q)-(q+2v)
    &=
    v^2+2v+q-q-2v
    =
    v^2,
    \label{eq:fator1_potts}
    \\[4pt]
    (v^2+2v+q)+(q+2v)
    &=
    v^2+2v+q+q+2v
    =
    v^2+4v+2q.
    \label{eq:fator2_potts}
\end{align}
%
Substituindo \eqref{eq:fator1_potts} e \eqref{eq:fator2_potts} em
\eqref{eq:diff_squares_potts},
%
\begin{equation}
    (v^2+2v+q)^2-(q+2v)^2
    =
    v^2(v^2+4v+2q).
    \label{eq:num_factored_potts}
\end{equation}
%
Portanto, a recorrência para \(v\) é
%
\begin{equation}
    \boxed{
    v'
    =
    \frac{
        v^2(v^2+4v+2q)
    }{(q+2v)^2}.
    }
    \label{eq:v_potts}
\end{equation}
% 
Os pontos fixos satisfazem
%
\begin{equation}
    v'=v.
\end{equation}
%
Da equação \eqref{eq:v_potts},
%
\begin{equation}
    v
    =
    \frac{
        v^2(v^2+4v+2q)
    }{(q+2v)^2}.
    \label{eq:fixed_v_potts}
\end{equation}
%
Existe imediatamente a solução
%
\begin{equation}
    \boxed{
    v^*=0,
    }
    \label{eq:v0_potts}
\end{equation}
%
que corresponde ao limite de alta temperatura, \(K\to0\). Para \(v\neq0\), podemos dividir ambos os lados de \eqref{eq:fixed_v_potts}
por \(v\):
%
\begin{equation}
    1
    =
    \frac{
        v(v^2+4v+2q)
    }{(q+2v)^2}.
\end{equation}
%
Multiplicando por \((q+2v)^2\),
%
\begin{equation}
    (q+2v)^2
    =
    v(v^2+4v+2q).
    \label{eq:fixed_v2_potts}
\end{equation}
%
Expandimos o lado esquerdo:
%
\begin{equation}
    (q+2v)^2
    =
    q^2+4qv+4v^2.
    \label{eq:lhs_exp_potts}
\end{equation}
%
Expandimos o lado direito:
%
\begin{equation}
    v(v^2+4v+2q)
    =
    v^3+4v^2+2qv.
    \label{eq:rhs_exp_potts}
\end{equation}
%
Igualando \eqref{eq:lhs_exp_potts} e \eqref{eq:rhs_exp_potts},
%
\begin{equation}
    q^2+4qv+4v^2
    =
    v^3+4v^2+2qv.
\end{equation}
%
Cancelando \(4v^2\) em ambos os lados,
%
\begin{equation}
    q^2+4qv
    =
    v^3+2qv.
\end{equation}
%
Rearranjando,
%
\begin{equation}
    v^3+2qv-4qv-q^2=0,
\end{equation}
%
ou seja,
%
\begin{equation}
    v^3-2qv-q^2=0.
    \label{eq:vc_potts}
\end{equation}
%
Essa equação cúbica determina o ponto crítico \(v_c\) para cada valor de \(q\). A equação \eqref{eq:vc_potts} é uma cúbica deprimida da forma
%
\begin{equation}
    v^3+pv+r=0,
    \qquad
    p=-2q,
    \qquad
    r=-q^2.
\end{equation}
%
Pela fórmula de Cardano, as soluções são
%
\begin{equation}
    v
    =
    \sqrt[3]{
        -\frac{r}{2}
        +
        \sqrt{
            \left(\frac{r}{2}\right)^2
            +
            \left(\frac{p}{3}\right)^3
        }
    }
    +
    \sqrt[3]{
        -\frac{r}{2}
        -
        \sqrt{
            \left(\frac{r}{2}\right)^2
            +
            \left(\frac{p}{3}\right)^3
        }
    }.
    \label{eq:cardano_general_potts}
\end{equation}
%
Calculamos os termos:
%
\begin{align}
    -\frac{r}{2}
    &=
    \frac{q^2}{2},
    \label{eq:cardano_term1_potts}
    \\[4pt]
    \left(\frac{r}{2}\right)^2
    &=
    \frac{q^4}{4},
    \label{eq:cardano_term2_potts}
    \\[4pt]
    \left(\frac{p}{3}\right)^3
    &=
    \left(-\frac{2q}{3}\right)^3
    =
    -\frac{8q^3}{27}.
    \label{eq:cardano_term3_potts}
\end{align}
%
Substituindo em \eqref{eq:cardano_general_potts},
%
\begin{equation}
    v_c
    =
    \sqrt[3]{
        \frac{q^2}{2}
        +
        \sqrt{
            \frac{q^4}{4}
            -
            \frac{8q^3}{27}
        }
    }
    +
    \sqrt[3]{
        \frac{q^2}{2}
        -
        \sqrt{
            \frac{q^4}{4}
            -
            \frac{8q^3}{27}
        }
    }.
    \label{eq:vc_cardano}
\end{equation}
%
Como
%
\begin{equation}
    K_c=\ln(1+v_c),
\end{equation}
%
a temperatura crítica é
%
\begin{equation}
    \boxed{
    \frac{k_BT_c(q)}{J}
    =
    \frac{1}{\ln[1+v_c(q)]}.
    }
    \label{eq:Tc_q}
\end{equation}
% 
Para obter o expoente térmico, linearizamos a transformação
%
\begin{equation}
    v'=R(v)
\end{equation}
%
em torno do ponto fixo \(v_c\):
%
\begin{equation}
    \delta v'
    =
    \lambda_t\,\delta v,
    \qquad
    \lambda_t=R'(v_c).
    \label{eq:linearization_potts}
\end{equation}
%
A partir de \eqref{eq:v_potts},
%
\begin{equation}
    R(v)
    =
    \frac{
        v^2(v^2+4v+2q)
    }{(q+2v)^2}.
    \label{eq:Rv_potts}
\end{equation}
%
Calculamos a derivada logarítmica de \(R(v)\):
%
\begin{equation}
    \ln R(v)
    =
    2\ln v
    +
    \ln(v^2+4v+2q)
    -
    2\ln(q+2v).
    \label{eq:lnR_potts}
\end{equation}
%
Derivando em relação a \(v\),
%
\begin{align}
    \frac{R'(v)}{R(v)}
    &=
    \frac{d}{dv}
    \left[
        2\ln v
        +
        \ln(v^2+4v+2q)
        -
        2\ln(q+2v)
    \right]
    \nonumber\\[4pt]
    &=
    \frac{2}{v}
    +
    \frac{2v+4}{v^2+4v+2q}
    -
    \frac{4}{q+2v}.
    \label{eq:logderiv_potts}
\end{align}
%
Como no ponto fixo \(R(v_c)=v_c\), podemos escrever
%
\begin{align}
    \lambda_t
    &=
    R'(v_c)
    =
    v_c\,
    \frac{R'(v_c)}{R(v_c)}
    \nonumber\\[4pt]
    &=
    v_c
    \left[
        \frac{2}{v_c}
        +
        \frac{2v_c+4}{v_c^2+4v_c+2q}
        -
        \frac{4}{q+2v_c}
    \right]
    \nonumber\\[4pt]
    &=
    2
    +
    \frac{
        2v_c(v_c+2)
    }{
        v_c^2+4v_c+2q
    }
    -
    \frac{
        4v_c
    }{
        q+2v_c
    }.
    \label{eq:lambdat_potts}
\end{align}
%
Como \(b=2\), o expoente \(y_t\) é definido por
%
\begin{equation}
    \lambda_t=b^{y_t}=2^{y_t},
\end{equation}
%
de modo que
%
\begin{equation}
    y_t(q)
    =
    \frac{\ln\lambda_t(q)}{\ln 2}.
    \label{eq:yt_potts}
\end{equation}
%
O expoente \(\nu\) é
%
\begin{equation}
    \nu(q)=\frac{1}{y_t(q)}.
    \label{eq:nu_potts}
\end{equation}
%
Usando a relação de hiperescala em \(d=2\),
%
\begin{equation}
    2-\alpha=2\nu,
\end{equation}
%
obtemos
%
\begin{equation}
    \boxed{
    \alpha(q)
    =
    2-\frac{2}{y_t(q)}
    =
    2-\frac{2\ln 2}{\ln\lambda_t(q)}.
    }
    \label{eq:alpha_q}
\end{equation}
% 
Alguns resultados numéricos são apresentados na
Tabela~\ref{tab:potts_values}.
%
\begin{table}[H]
    \centering
    \caption{
        Temperatura crítica, expoente \(\nu\) e expoente \(\alpha\)
        para alguns valores de \(q\) na rede hierárquica do diamante.
    }
    \label{tab:potts_values}
    \begin{tabular}{cccc}
        \toprule
        \(q\) & \(K_c\) & \(\nu\) & \(\alpha\) \\
        \midrule
        2 & 1.21876 & 1.33827 & $-0.67653$ \\
        3 & 1.38629 & 1.20471 & $-0.40942$ \\
        4 & 1.51262 & 1.12416 & $-0.24832$ \\
        5 & 1.61464 & 1.06865 & $-0.13729$ \\
        6 & 1.70052 & 1.02731 & $-0.05463$ \\
        7 & 1.77481 & 0.99494 & $0.01012$ \\
        8 & 1.84037 & 0.96866 & $0.06268$ \\
        9 & 1.89910 & 0.94675 & $0.10649$ \\
        10 & 1.95232 & 0.92811 & $0.14378$ \\
        \bottomrule
    \end{tabular}
\end{table}
%
Para \(q=2\), o modelo de Potts é equivalente ao modelo de Ising.
De fato, com \(\sigma_i=\pm1\),
%
\begin{equation}
    \delta_{\sigma_i,\sigma_j}
    =
    \frac{1+\sigma_i\sigma_j}{2},
\end{equation}
%
de modo que
%
\begin{equation}
    -J\sum_{\langle ij\rangle}\delta_{\sigma_i,\sigma_j}
    =
    -\frac{J}{2}\sum_{\langle ij\rangle}\sigma_i\sigma_j
    -
    \frac{J}{2}\sum_{\langle ij\rangle}1.
\end{equation}
%
O último termo é uma constante, e o acoplamento efetivo de Ising é
\(J_{\mathrm{Ising}}=J/2\). Assim, \(K_{\mathrm{Ising}}=K_{\mathrm{Potts}}/2\).
Para \(q=2\), a equação \eqref{eq:vc_potts} dá
%
\begin{equation}
    v_c^3-4v_c-4=0,
\end{equation}
%
cuja solução positiva é \(v_c\simeq2.38298\). Então
%
\begin{equation}
    K_c^{\mathrm{Potts}}
    =
    \ln(1+v_c)
    \simeq
    1.21876,
\end{equation}
%
e
%
\begin{equation}
    K_c^{\mathrm{Ising}}
    =
    \frac{K_c^{\mathrm{Potts}}}{2}
    \simeq
    0.60938,
\end{equation}
%
que coincide exatamente com o valor obtido no Problema~1. O expoente \(\alpha(q)\) passa de valores negativos para positivos
em torno de
%
\begin{equation}
    q\simeq6.83.
\end{equation}
%
A dependência de \(\alpha\) com \(q\) é mostrada na
Fig.~\ref{fig:alpha_potts}.
%
\begin{figure}[H]
    \centering
    \includegraphics[width=0.82\linewidth]{Figures/statistical_mechanics/alpha_potts_diamond_mk.pdf}
    \caption{
        Expoente térmico \(\alpha\) do modelo de Potts na rede
        hierárquica do diamante, obtido a partir da linearização da
        transformação de renormalização e da relação de hiperescala
        em \(d=2\).
    }
    \label{fig:alpha_potts}
\end{figure}
% ============================================================
% PROBLEMA 3
% ============================================================
\section*{Problema 3 -- Renormalização em blocos de cinco spins}
%
Neste problema utilizamos uma transformação de grupo de renormalização no espaço real baseada em aglomerados de cinco spins em forma de cruz, seguindo a geometria indicada na figura da lista. O Hamiltoniano de
partida é
%
\begin{equation}
    \mathcal H
    =
    -J\sum_{\langle ij\rangle}
    \sigma_i\sigma_j,
    \qquad
    \sigma_i=\pm1,
    \qquad
    J>0.
    \label{eq:H_p3}
\end{equation}
%
Cada bloco contém cinco spins: um spin central \(\sigma_0\) e quatro spins de borda \(\sigma_1,\sigma_2,\sigma_3,\sigma_4\). A variável renormalizada \(\Sigma\) é definida pela regra da maioria:
%
\begin{equation}
    \Sigma
    =
    \operatorname{sgn}
    \left(
        \sigma_0+\sigma_1+\sigma_2+\sigma_3+\sigma_4
    \right).
    \label{eq:majority}
\end{equation}
%
Como há cinco spins, a soma nunca se anula e, portanto, não existe empate. A nova rede é uma rede quadrada girada. Para ver isso, note que os centros dos blocos podem ser tomados, por exemplo, nos pontos \((2m-n,\,m+2n)\), com \(m,n\in\mathbb Z\). Dois centros vizinhos,
como \((0,0)\) e \((2,1)\), estão separados por
%
\begin{equation}
    \sqrt{2^2+1^2}=\sqrt5.
\end{equation}
%
Logo, o parâmetro de rede é multiplicado por
%
\begin{equation}
    b=\sqrt5.
    \label{eq:b_p3}
\end{equation}
%
Denotemos por \(\sigma_0\) o spin central e por
\(\sigma_1,\ldots,\sigma_4\) os quatro spins da borda. O Hamiltoniano
interno do bloco é
%
\begin{equation}
    \mathcal H_B
    =
    -J\sigma_0
    \sum_{a=1}^{4}\sigma_a.
\end{equation}
%
Introduzindo \(K=\beta J\), a probabilidade interna é proporcional a
%
\begin{equation}
    \exp\left[
        K\sigma_0\sum_{a=1}^{4}\sigma_a
    \right].
\end{equation}
%
Para calcular a interação efetiva entre blocos, precisamos da média
de um spin de borda condicionada à variável de bloco \(\Sigma\).
Por simetria, todas as quatro bordas possuem a mesma média. Definimos
%
\begin{equation}
    \mu(K)
    \equiv
    \left\langle
        \sigma_1
    \right\rangle_{\Sigma=+1}.
\end{equation}
%
Para calcular \(\mu(K)\), escrevemos
%
\begin{equation}
    S\equiv \sigma_1+\sigma_2+\sigma_3+\sigma_4.
\end{equation}
%
Se \(k\) é o número de spins \(+1\) entre as quatro bordas, então
%
\begin{equation}
    S=2k-4,
    \qquad
    k=0,1,2,3,4.
\end{equation}
%
A condição de maioria positiva é
%
\begin{equation}
    \Sigma=+1
    \quad\Longleftrightarrow\quad
    \sigma_0+S>0.
\end{equation}
%
Analisamos separadamente \(\sigma_0=+1\) e \(\sigma_0=-1\).
%
\begin{itemize}
    \item Se \(\sigma_0=+1\), a condição \(1+S>0\) exige \(S\ge0\), ou seja,
    \(S=0,2,4\), que correspondem a \(k=2,3,4\).

    \item Se \(\sigma_0=-1\), a condição \(-1+S>0\) exige \(S\ge2\), ou seja,
    \(S=2,4\), que correspondem a \(k=3,4\).
\end{itemize}
%
O denominador da média, que chamaremos de \(Z_+\), é a soma dos pesos de Boltzmann sobre todas as configurações com maioria positiva:
%
\begin{align}
    Z_+
    &=
    \sum_{\substack{\text{configs}\\\Sigma=+1}}
    e^{K\sigma_0 S}
    \nonumber\\[4pt]
    &=
    \underbrace{
        \binom{4}{2}\,e^{K(+1)(0)}
    }_{k=2,\ \sigma_0=+1}
    +
    \underbrace{
        \binom{4}{3}\,e^{K(+1)(2)}
    }_{k=3,\ \sigma_0=+1}
    +
    \underbrace{
        \binom{4}{4}\,e^{K(+1)(4)}
    }_{k=4,\ \sigma_0=+1}
    \nonumber\\[4pt]
    &\qquad
    +
    \underbrace{
        \binom{4}{3}\,e^{K(-1)(2)}
    }_{k=3,\ \sigma_0=-1}
    +
    \underbrace{
        \binom{4}{4}\,e^{K(-1)(4)}
    }_{k=4,\ \sigma_0=-1}.
\end{align}
%
Como
%
\begin{equation}
    \binom{4}{2}=6,
    \qquad
    \binom{4}{3}=4,
    \qquad
    \binom{4}{4}=1,
\end{equation}
%
segue
%
\begin{equation}
    Z_+
    =
    6
    +
    4e^{2K}
    +
    e^{4K}
    +
    4e^{-2K}
    +
    e^{-4K}.
    \label{eq:Zplus_p3}
\end{equation}
%
Multiplicando \eqref{eq:Zplus_p3} por \(e^{4K}\), obtemos
%
\begin{align}
    e^{4K}Z_+
    &=
    6e^{4K}
    +
    4e^{6K}
    +
    e^{8K}
    +
    4e^{2K}
    +
    1
    \nonumber\\[4pt]
    &=
    1+4e^{2K}+6e^{4K}+4e^{6K}+e^{8K}
    \nonumber\\[4pt]
    &=
    \left(1+e^{2K}\right)^4.
    \label{eq:Zplus_factored_p3}
\end{align}
%
Passemos agora ao numerador \(N_+\), que contém o fator \(\sigma_1\):
%
\begin{equation}
    N_+
    =
    \sum_{\substack{\text{configs}\\\Sigma=+1}}
    \sigma_1\,
    e^{K\sigma_0 S}.
\end{equation}
%
Para uma dada borda \(\sigma_1\), a soma sobre as demais bordas deve
ser feita com cuidado. Fixado \(k\), o número de configurações com
\(\sigma_1=+1\) é \(\binom{3}{k-1}\), e o número com \(\sigma_1=-1\)
é \(\binom{3}{k}\). Portanto, a soma de \(\sigma_1\) sobre as
configurações de borda com um dado \(k\) é
%
\begin{equation}
    \sum_{\text{bordas}}\sigma_1
    =
    \binom{3}{k-1}
    -
    \binom{3}{k}.
    \label{eq:sum_sigma1_p3}
\end{equation}
%
Calculamos \eqref{eq:sum_sigma1_p3} para os valores de \(k\) que
aparecem em \(Z_+\):
%
\begin{align}
    k=2:&\qquad
    \binom{3}{1}-\binom{3}{2}
    =
    3-3
    =
    0,
    \nonumber\\[4pt]
    k=3:&\qquad
    \binom{3}{2}-\binom{3}{3}
    =
    3-1
    =
    2,
    \nonumber\\[4pt]
    k=4:&\qquad
    \binom{3}{3}-\binom{3}{4}
    =
    1-0
    =
    1.
\end{align}
%
Assim, as contribuições para \(N_+\) são:
%
\begin{align}
    N_+
    &=
    \underbrace{
        0\cdot e^{K(+1)(0)}
    }_{k=2,\ \sigma_0=+1}
    +
    \underbrace{
        2\,e^{K(+1)(2)}
    }_{k=3,\ \sigma_0=+1}
    +
    \underbrace{
        1\,e^{K(+1)(4)}
    }_{k=4,\ \sigma_0=+1}
    \nonumber\\[4pt]
    &\qquad
    +
    \underbrace{
        2\,e^{K(-1)(2)}
    }_{k=3,\ \sigma_0=-1}
    +
    \underbrace{
        1\,e^{K(-1)(4)}
    }_{k=4,\ \sigma_0=-1}
    \nonumber\\[4pt]
    &=
    2e^{2K}
    +
    e^{4K}
    +
    2e^{-2K}
    +
    e^{-4K}.
    \label{eq:Nplus_p3}
\end{align}
%
Multiplicando \eqref{eq:Nplus_p3} por \(e^{4K}\),
%
\begin{align}
    e^{4K}N_+
    &=
    2e^{6K}
    +
    e^{8K}
    +
    2e^{2K}
    +
    1
    \nonumber\\[4pt]
    &=
    e^{8K}+2e^{6K}+2e^{2K}+1.
    \label{eq:Nplus_factored_p3}
\end{align}
%
Dividindo \eqref{eq:Nplus_factored_p3} por \eqref{eq:Zplus_factored_p3},
obtemos
%
\begin{equation}
    \mu(K)
    =
    \frac{
        e^{8K}+2e^{6K}+2e^{2K}+1
    }
    {\left(1+e^{2K}\right)^4}.
    \label{eq:mu_exp_p3}
\end{equation}
% 
Introduzimos
%
\begin{equation}
    t=\tanh K.
\end{equation}
%
Queremos simplificar \eqref{eq:mu_exp_p3}. Seja
%
\begin{equation}
    u=e^{2K}.
\end{equation}
%
Então
%
\begin{equation}
    e^{8K}=u^4,
    \qquad
    e^{6K}=u^3,
    \qquad
    e^{2K}=u,
\end{equation}
%
e a expressão \eqref{eq:mu_exp_p3} torna-se
%
\begin{equation}
    \mu
    =
    \frac{
        u^4+2u^3+2u+1
    }
    {(1+u)^4}.
    \label{eq:mu_u_p3}
\end{equation}
%
Dividimos numerador e denominador de \eqref{eq:mu_u_p3} por \(u^2\):
%
\begin{equation}
    \mu
    =
    \frac{
        u^2+2u+\dfrac{2}{u}+\dfrac{1}{u^2}
    }
    {
        \dfrac{(1+u)^4}{u^2}
    }.
    \label{eq:mu_u_div_p3}
\end{equation}
%
No denominador,
%
\begin{equation}
    \frac{(1+u)^4}{u^2}
    =
    \left[
        \frac{(1+u)^2}{u}
    \right]^2
    =
    \left(
        u+2+\frac{1}{u}
    \right)^2.
\end{equation}
%
Definindo
%
\begin{equation}
    w\equiv u+\frac{1}{u},
\end{equation}
%
o denominador fica
%
\begin{equation}
    (w+2)^2.
\end{equation}
%
No numerador,
%
\begin{align}
    u^2+2u+\frac{2}{u}+\frac{1}{u^2}
    &=
    \left(
        u^2+2+\frac{1}{u^2}
    \right)
    +
    2u+\frac{2}{u}-2
    \nonumber\\[4pt]
    &=
    \left(
        u+\frac{1}{u}
    \right)^2
    +
    2\left(
        u+\frac{1}{u}
    \right)
    -2
    \nonumber\\[4pt]
    &=
    w^2+2w-2.
\end{align}
%
Portanto,
%
\begin{equation}
    \mu
    =
    \frac{w^2+2w-2}{(w+2)^2}.
    \label{eq:mu_w_p3}
\end{equation}
%
Agora usamos a relação entre \(u\) e \(t\). Como
%
\begin{equation}
    t=\tanh K
    =
    \frac{e^{2K}-1}{e^{2K}+1}
    =
    \frac{u-1}{u+1},
\end{equation}
%
podemos inverter:
%
\begin{equation}
    u=\frac{1+t}{1-t}.
    \label{eq:u_t_p3}
\end{equation}
%
Logo,
%
\begin{align}
    w
    &=
    u+\frac{1}{u}
    =
    \frac{1+t}{1-t}
    +
    \frac{1-t}{1+t}
    \nonumber\\[4pt]
    &=
    \frac{(1+t)^2+(1-t)^2}{(1-t)(1+t)}
    \nonumber\\[4pt]
    &=
    \frac{1+2t+t^2+1-2t+t^2}{1-t^2}
    \nonumber\\[4pt]
    &=
    \frac{2+2t^2}{1-t^2}
    =
    \frac{2(1+t^2)}{1-t^2}.
    \label{eq:w_t_p3}
\end{align}
%
Assim,
%
\begin{equation}
    w+2
    =
    \frac{2(1+t^2)}{1-t^2}+2
    =
    \frac{2(1+t^2)+2(1-t^2)}{1-t^2}
    =
    \frac{4}{1-t^2}.
    \label{eq:wplus2_t_p3}
\end{equation}
%
Além disso,
%
\begin{align}
    w^2+2w-2
    &=
    \left[
        \frac{2(1+t^2)}{1-t^2}
    \right]^2
    +
    2\left[
        \frac{2(1+t^2)}{1-t^2}
    \right]
    -2
    \nonumber\\[4pt]
    &=
    \frac{
        4(1+t^2)^2
        +
        4(1+t^2)(1-t^2)
        -
        2(1-t^2)^2
    }
    {(1-t^2)^2}.
    \label{eq:num_common_p3}
\end{align}
%
Calculamos cada termo do numerador de \eqref{eq:num_common_p3}:
%
\begin{align}
    4(1+t^2)^2
    &=
    4(1+2t^2+t^4)
    =
    4+8t^2+4t^4,
    \nonumber\\[4pt]
    4(1+t^2)(1-t^2)
    &=
    4(1-t^4)
    =
    4-4t^4,
    \nonumber\\[4pt]
    -2(1-t^2)^2
    &=
    -2(1-2t^2+t^4)
    =
    -2+4t^2-2t^4.
\end{align}
%
Somando,
%
\begin{equation}
    4+8t^2+4t^4
    +
    4-4t^4
    -
    2+4t^2-2t^4
    =
    6+12t^2-2t^4.
\end{equation}
%
Portanto,
%
\begin{equation}
    w^2+2w-2
    =
    \frac{
        6+12t^2-2t^4
    }
    {(1-t^2)^2}
    =
    \frac{
        2(3+6t^2-t^4)
    }
    {(1-t^2)^2}.
    \label{eq:num_final_p3}
\end{equation}
%
Substituindo \eqref{eq:num_final_p3} e \eqref{eq:wplus2_t_p3} em
\eqref{eq:mu_w_p3},
%
\begin{equation}
    \mu
    =
    \frac{
        \dfrac{2(3+6t^2-t^4)}{(1-t^2)^2}
    }
    {
        \left(
            \dfrac{4}{1-t^2}
        \right)^2
    }
    =
    \frac{
        \dfrac{2(3+6t^2-t^4)}{(1-t^2)^2}
    }
    {
        \dfrac{16}{(1-t^2)^2}
    }.
\end{equation}
%
Cancelando \((1-t^2)^2\),
%
\begin{equation}
    \mu(K)
    =
    \frac{
        3+6\tanh^2 K-\tanh^4 K
    }{8}.
    \label{eq:mu_p3}
\end{equation}
%
Por simetria global,
%
\begin{equation}
    \left\langle
        \sigma_a
    \right\rangle_{\Sigma}
    =
    \mu(K)\Sigma,
    \qquad
    a=1,2,3,4.
\end{equation}
% 
Pela geometria dos blocos, dois blocos vizinhos \(A\) e \(B\) são
conectados por dois bonds da rede quadrada original. Assim, a parte
do Hamiltoniano que acopla dois blocos pode ser escrita como
%
\begin{equation}
    \mathcal H_{AB}
    =
    -J
    \left(
        \sigma_a\sigma_b
        +
        \sigma_c\sigma_d
    \right),
    \label{eq:HAB_p3}
\end{equation}
%
onde \(\sigma_a,\sigma_c\) pertencem ao bloco \(A\) e
\(\sigma_b,\sigma_d\) pertencem ao bloco \(B\). No cálculo perturbativo até primeira ordem em \(K=\beta J\),
%
\begin{equation}
    e^{-\beta\mathcal H_{AB}}
    \simeq
    1-\beta\mathcal H_{AB}
    =
    1
    +
    K
    \left(
        \sigma_a\sigma_b
        +
        \sigma_c\sigma_d
    \right).
    \label{eq:exp_pert_p3}
\end{equation}
%
Tomando a média sobre os estados internos dos dois blocos, e usando
que os blocos são independentes,
%
\begin{align}
    \left\langle
        \sigma_a\sigma_b
    \right\rangle
    &=
    \left\langle
        \sigma_a
    \right\rangle_A
    \left\langle
        \sigma_b
    \right\rangle_B
    =
    \mu(K)\Sigma_A\,
    \mu(K)\Sigma_B
    \nonumber\\[4pt]
    &=
    \mu(K)^2\,\Sigma_A\Sigma_B.
\end{align}
%
Analogamente,
%
\begin{equation}
    \left\langle
        \sigma_c\sigma_d
    \right\rangle
    =
    \mu(K)^2\,\Sigma_A\Sigma_B.
\end{equation}
%
Portanto,
%
\begin{align}
    -\beta
    \left\langle
        \mathcal H_{AB}
    \right\rangle
    &=
    K
    \left[
        \left\langle
            \sigma_a\sigma_b
        \right\rangle
        +
        \left\langle
            \sigma_c\sigma_d
        \right\rangle
    \right]
    \nonumber\\[4pt]
    &=
    K
    \left[
        \mu^2\Sigma_A\Sigma_B
        +
        \mu^2\Sigma_A\Sigma_B
    \right]
    \nonumber\\[4pt]
    &=
    2K\mu(K)^2\,\Sigma_A\Sigma_B.
\end{align}
%
Logo, a interação efetiva entre os spins renormalizados é
%
\begin{equation}
    -K'\Sigma_A\Sigma_B,
\end{equation}
%
com
%
\begin{equation}
    K'
    =
    2K\mu(K)^2.
    \label{eq:RG_cross}
\end{equation}
%
Substituindo \eqref{eq:mu_p3},
%
\begin{equation}
    \boxed{
    K'
    =
    \frac{K}{32}
    \left[
        3+6\tanh^2K-\tanh^4K
    \right]^2.
    }
    \label{eq:RG_cross_explicit}
\end{equation}
%
O ponto crítico satisfaz
%
\begin{equation}
    K_c=2K_c\mu(K_c)^2.
\end{equation}
%
Para \(K_c\neq0\), podemos dividir por \(K_c\):
%
\begin{equation}
    1=2\mu(K_c)^2,
\end{equation}
%
ou seja,
%
\begin{equation}
    \mu(K_c)=\frac{1}{\sqrt2}.
    \label{eq:mu_c_p3}
\end{equation}
%
Usando \eqref{eq:mu_p3},
%
\begin{equation}
    \frac{
        3+6\tanh^2K_c-\tanh^4K_c
    }{8}
    =
    \frac{1}{\sqrt2}.
    \label{eq:eq_Kc_p3}
\end{equation}
%
Multiplicando por \(8\),
%
\begin{equation}
    3+6\tanh^2K_c-\tanh^4K_c
    =
    \frac{8}{\sqrt2}
    =
    4\sqrt2.
\end{equation}
%
Rearranjando,
%
\begin{equation}
    \tanh^4K_c
    -
    6\tanh^2K_c
    +
    (4\sqrt2-3)
    =
    0.
    \label{eq:quartic_t_p3}
\end{equation}
%
Definindo
%
\begin{equation}
    y\equiv \tanh^2 K_c,
\end{equation}
%
a equação \eqref{eq:quartic_t_p3} torna-se uma equação quadrática:
%
\begin{equation}
    y^2-6y+(4\sqrt2-3)=0.
\end{equation}
%
As soluções são
%
\begin{align}
    y
    &=
    \frac{
        6\pm\sqrt{36-4(4\sqrt2-3)}
    }{2}
    \nonumber\\[4pt]
    &=
    3\pm\sqrt{9-(4\sqrt2-3)}
    \nonumber\\[4pt]
    &=
    3\pm\sqrt{12-4\sqrt2}
    \nonumber\\[4pt]
    &=
    3\pm 2\sqrt{3-\sqrt2}.
\end{align}
%
Como \(0<y<1\), escolhemos o sinal negativo:
%
\begin{equation}
    \tanh^2 K_c
    =
    3-2\sqrt{3-\sqrt2}
    \simeq
    0.4814.
    \label{eq:tanh2_Kc_p3}
\end{equation}
%
Portanto,
%
\begin{equation}
    \tanh K_c
    =
    \sqrt{
        3-2\sqrt{3-\sqrt2}
    }
    \simeq
    0.6938.
\end{equation}
%
O acoplamento crítico é
%
\begin{equation}
    K_c
    =
    \operatorname{arctanh}
    \left(
        \tanh K_c
    \right)
    =
    \frac12
    \ln
    \left(
        \frac{
            1+\tanh K_c
        }{
            1-\tanh K_c
        }
    \right)
    \simeq
    0.855359.
    \label{eq:Kc_p3}
\end{equation}
%
Assim,
%
\begin{equation}
    \boxed{
    \frac{k_BT_c}{J}
    =
    \frac1{K_c}
    \simeq
    1.16910.
    }
    \label{eq:Tc_p3}
\end{equation}
% 
Linearizando a transformação
%
\begin{equation}
    K'=R(K)=2K\mu(K)^2
\end{equation}
%
em torno de \(K_c\),
%
\begin{equation}
    \delta K'
    =
    R'(K_c)\,\delta K,
    \qquad
    \lambda_t
    =
    R'(K_c).
\end{equation}
%
Derivando,
%
\begin{equation}
    \lambda_t
    =
    2\mu(K_c)^2
    +
    4K_c\mu(K_c)\mu'(K_c).
\end{equation}
%
Como \(\mu(K_c)^2=1/2\),
%
\begin{equation}
    \lambda_t
    =
    1
    +
    4K_c\mu(K_c)\mu'(K_c).
    \label{eq:lambda_form_p3}
\end{equation}
%
Precisamos de \(\mu'(K)\). A partir de \eqref{eq:mu_p3},
%
\begin{equation}
    \mu(K)
    =
    \frac{
        3+6t^2-t^4
    }{8},
    \qquad
    t=\tanh K.
\end{equation}
%
Derivando em relação a \(K\),
%
\begin{equation}
    \frac{dt}{dK}
    =
    \operatorname{sech}^2K
    =
    1-\tanh^2K
    =
    1-t^2.
\end{equation}
%
Logo,
%
\begin{align}
    \mu'(K)
    &=
    \frac18
    \left(
        12t-4t^3
    \right)
    \frac{dt}{dK}
    \nonumber\\[4pt]
    &=
    \frac18
    \left(
        12t-4t^3
    \right)
    (1-t^2)
    \nonumber\\[4pt]
    &=
    \frac{4t(3-t^2)}{8}
    (1-t^2)
    \nonumber\\[4pt]
    &=
    \frac{t}{2}
    (3-t^2)(1-t^2).
    \label{eq:mu_prime_p3}
\end{align}
%
Usando os valores numéricos
%
\begin{equation}
    t_c=\tanh K_c\simeq0.6938,
    \qquad
    \mu_c=\frac{1}{\sqrt2}\simeq0.7071,
    \qquad
    K_c\simeq0.855359,
\end{equation}
%
e
%
\begin{equation}
    \mu'(K_c)
    \simeq
    0.4528,
\end{equation}
%
obtemos
%
\begin{equation}
    \lambda_t
    \simeq
    1
    +
    4(0.855359)(0.7071)(0.4528)
    \simeq
    2.09619.
    \label{eq:lambda_p3}
\end{equation}
%
Como \(b=\sqrt5\),
%
\begin{equation}
    \lambda_t=b^{y_t},
\end{equation}
%
e, portanto,
%
\begin{equation}
    y_t
    =
    \frac{\ln\lambda_t}{\ln b}
    =
    \frac{\ln\lambda_t}{\ln\sqrt5}
    \simeq
    0.91973.
\end{equation}
%
Logo,
%
\begin{equation}
    \boxed{
    \nu
    =
    \frac1{y_t}
    \simeq
    1.08728.
    }
    \label{eq:nu_p3}
\end{equation}
%
Aplicando a relação de hiperescala em \(d=2\),
%
\begin{equation}
    \alpha
    =
    2-2\nu,
\end{equation}
%
obtemos
%
\begin{equation}
    \boxed{
    \alpha
    \simeq
    -0.17456.
    }
    \label{eq:alpha_p3}
\end{equation}
% 
Esse resultado melhora, dentro dessa aproximação de primeira ordem, em relação ao valor obtido pela construção do diamante,
\(\alpha_{\mathrm{MK}}\simeq-0.6765\), embora ainda esteja distante
do resultado exato \(\alpha=0\) para o modelo de Ising bidimensional. A diferença residual deve-se ao caráter perturbativo de primeira ordem e ao fato de que a transformação de renormalização é truncada, não incorporando todas as correlações geradas pelos blocos.
% PROBLEMA 4
% ============================================================
\section*{Problema 4 -- Modelo de Ising de spin 1 e geração de novas interações}
%
Consideramos inicialmente o modelo de Ising de spin 1
%
\begin{equation}
    \mathcal H
    =
    -J\sum_{\langle ij\rangle}S_iS_j,
    \qquad
    S_i=-1,0,+1,
    \qquad
    J>0.
    \label{eq:H_spin1_initial}
\end{equation}
%
O objetivo é mostrar que uma transformação de renormalização não pode ser fechada apenas em termos da constante \(J\) e, posteriormente, estudar o espaço ampliado de Hamiltonianas contendo interações adicionais. Consideremos uma cadeia
%
\begin{equation}
    S_1-S-S_2
\end{equation}
%
e eliminemos o spin intermediário \(S\). O peso efetivo é
%
\begin{equation}
    W(S_1,S_2)
    =
    \sum_{S=-1}^{1}
    e^{K S(S_1+S_2)},
    \qquad
    K=\beta J.
    \label{eq:Wspin1}
\end{equation}
%
Calculamos explicitamente \(W(S_1,S_2)\) para todos os casos distintos. Para o caso \(S_1=S_2=+1\), \(S_1+S_2=2\). Então
%
\begin{align}
    W(+1,+1)
    &=
    \sum_{S=-1}^{1}
    e^{2K S}
    \nonumber\\[4pt]
    &=
    e^{-2K}
    +
    e^{0}
    +
    e^{2K}
    \nonumber\\[4pt]
    &=
    1
    +
    2\cosh(2K).
    \label{eq:W++_p4}
\end{align}
%
Para o caso \(S_1=+1\), \(S_2=-1\), \(S_1+S_2=0\). Logo,
%
\begin{equation}
    W(+1,-1)
    =
    \sum_{S=-1}^{1}
    e^{0}
    =
    1+1+1
    =
    3.
    \label{eq:W+-_p4}
\end{equation}
%
Para o caso \(S_1=+1\), \(S_2=0\), \(S_1+S_2=1\). Então
%
\begin{align}
    W(+1,0)
    &=
    \sum_{S=-1}^{1}
    e^{K S}
    \nonumber\\[4pt]
    &=
    e^{-K}
    +
    1
    +
    e^{K}
    \nonumber\\[4pt]
    &=
    1
    +
    2\cosh K.
    \label{eq:W+0_p4}
\end{align}
%
Para o caso \(S_1=S_2=0\), \(S_1+S_2=0\). Portanto,
%
\begin{equation}
    W(0,0)
    =
    \sum_{S=-1}^{1}
    e^{0}
    =
    3.
    \label{eq:W00_p4}
\end{equation}
%
Suponha que a transformação pudesse ser escrita novamente apenas como uma interação bilinear,
%
\begin{equation}
    \mathcal H'
    =
    -J'S_1S_2
\end{equation}
%
mais uma constante. Nesse caso, o peso efetivo teria a forma
%
\begin{equation}
    W_{\mathrm{eff}}(S_1,S_2)
    =
    A\,e^{K' S_1S_2},
    \qquad
    K'=\beta J'.
\end{equation}
%
Para \(S_1=+1\) e \(S_2=0\), temos \(S_1S_2=0\). Para \(S_1=S_2=0\),
também temos \(S_1S_2=0\). Logo,
%
\begin{equation}
    W_{\mathrm{eff}}(+1,0)
    =
    A\,e^{0}
    =
    A,
    \qquad
    W_{\mathrm{eff}}(0,0)
    =
    A\,e^{0}
    =
    A.
\end{equation}
%
Portanto, deveríamos ter necessariamente
%
\begin{equation}
    W(+1,0)=W(0,0).
\end{equation}
%
Entretanto, de \eqref{eq:W+0_p4} e \eqref{eq:W00_p4},
%
\begin{equation}
    W(+1,0)=1+2\cosh K,
    \qquad
    W(0,0)=3.
\end{equation}
%
Para \(K\neq0\),
%
\begin{equation}
    1+2\cosh K
    \neq
    3,
\end{equation}
%
pois \(\cosh K>1\) para \(K\neq0\). Logo, a transformação não fecha no espaço de Hamiltonianas parametrizado apenas por $J$. Novas interações são geradas pela decimação. Para obter um espaço de Hamiltonianas fechado, introduzimos
%
\begin{equation}
    \mathcal H
    =
    -J\sum_{\langle ij\rangle}S_iS_j
    -K\sum_{\langle ij\rangle}S_i^2S_j^2
    -D\sum_iS_i^2.
    \label{eq:H_spin1_extended}
\end{equation}
%
O primeiro termo é a interação bilinear, o segundo é a interação
biquadrática e o terceiro é o campo cristalino. Seguindo o enunciado, introduzimos as variáveis adimensionais
%
\begin{equation}
    x=
    \exp\left[
        -\frac{\beta}{2}(J+K+D)
    \right],
    \label{eq:x_spin1}
\end{equation}
%
\begin{equation}
    y=e^{-2\beta J},
    \label{eq:y_spin1}
\end{equation}
%
e
%
\begin{equation}
    z=e^{-\beta(J+K+D)}.
    \label{eq:z_spin1}
\end{equation}
%
Observamos imediatamente que, para os parâmetros originais,
%
\begin{equation}
    z
    =
    e^{-\beta(J+K+D)}
    =
    \left[
        e^{-\frac{\beta}{2}(J+K+D)}
    \right]^2
    =
    x^2.
    \label{eq:z_x2_initial}
\end{equation}
%
No entanto, a transformação de renormalização gera, em geral, três variáveis independentes \(x',y',z'\). Esse é precisamente o fechamento do espaço de Hamiltonianas que faltava no modelo com apenas \(J\). Ao eliminar os graus de liberdade intermediários e normalizar os
pesos efetivos, obtêm-se as recorrências fornecidas no enunciado:
%
\begin{equation}
    x'
    =
    \frac{(1+y+z)x}
         {1+x^2+y^2},
    \label{eq:xrec}
\end{equation}
%
\begin{equation}
    y'
    =
    \frac{2y+x^2}
         {1+x^2+y^2},
    \label{eq:yrec}
\end{equation}
%
e
%
\begin{equation}
    z'
    =
    \frac{z^2+2x^2}
         {1+x^2+y^2}.
    \label{eq:zrec}
\end{equation}
%
Para entender a estrutura dessas equações, definimos o denominador
comum
%
\begin{equation}
    \mathcal D=1+x^2+y^2.
    \label{eq:D_p4}
\end{equation}
%
Os pesos renormalizados correspondentes aos três setores são
%
\begin{equation}
    \widetilde x=(1+y+z)x,
    \label{eq:x_tilde_p4}
\end{equation}
%
\begin{equation}
    \widetilde y=2y+x^2,
    \label{eq:y_tilde_p4}
\end{equation}
%
e
%
\begin{equation}
    \widetilde z=z^2+2x^2.
    \label{eq:z_tilde_p4}
\end{equation}
%
Normalizando pelo peso de referência \(\mathcal D\),
%
\begin{equation}
    x'=\frac{\widetilde x}{\mathcal D},
    \qquad
    y'=\frac{\widetilde y}{\mathcal D},
    \qquad
    z'=\frac{\widetilde z}{\mathcal D}.
\end{equation}
%
Isso reproduz \eqref{eq:xrec}--\eqref{eq:zrec}. Os pontos fixos satisfazem simultaneamente
%
\begin{equation}
    x'=x,
    \qquad
    y'=y,
    \qquad
    z'=z.
    \label{eq:fixed_conditions_p4}
\end{equation}
%
Agora, par o subcaso \(x>0\). A condição \(x'=x\) implica
%
\begin{equation}
    x
    =
    \frac{(1+y+z)x}{\mathcal D}.
\end{equation}
%
Como \(x>0\), podemos dividir ambos os lados por \(x\):
%
\begin{equation}
    1
    =
    \frac{1+y+z}{\mathcal D},
\end{equation}
%
ou seja,
%
\begin{equation}
    \mathcal D=1+y+z.
    \label{eq:D_eq_p4}
\end{equation}
%
Substituindo a definição de \(\mathcal D\),
%
\begin{equation}
    1+x^2+y^2
    =
    1+y+z.
\end{equation}
%
Cancelando \(1\) em ambos os lados,
%
\begin{equation}
    x^2+y^2
    =
    y+z,
\end{equation}
%
ou
%
\begin{equation}
    x^2+y^2-y-z=0.
    \label{eq:fp1}
\end{equation}

A condição \(y'=y\) fornece
%
\begin{equation}
    y
    =
    \frac{2y+x^2}{\mathcal D}.
\end{equation}
%
Multiplicando por \(\mathcal D\),
%
\begin{equation}
    y(1+x^2+y^2)
    =
    2y+x^2.
\end{equation}
%
Expandindo o lado esquerdo,
%
\begin{equation}
    y+x^2y+y^3
    =
    2y+x^2.
\end{equation}
%
Rearranjando,
%
\begin{equation}
    x^2y+y^3-y-x^2=0.
\end{equation}
%
Fatorando,
%
\begin{align}
    x^2y+y^3-y-x^2
    &=
    x^2(y-1)+y(y^2-1)
    \nonumber\\[4pt]
    &=
    x^2(y-1)+y(y-1)(y+1)
    \nonumber\\[4pt]
    &=
    (y-1)\left[
        x^2+y(y+1)
    \right].
\end{align}
%
Portanto,
%
\begin{equation}
    (y-1)
    \left[
        x^2+y(y+1)
    \right]
    =0.
    \label{eq:fp2}
\end{equation}
%
No domínio físico \(x,y\ge0\), o segundo fator é não negativo, pois
\(x^2\ge0\) e \(y(y+1)\ge0\). Ele só se anula se \(x=0\) e \(y=0\).
Como estamos no caso \(x>0\), o segundo fator é estritamente positivo.
Assim, devemos ter
%
\begin{equation}
    y=1.
    \label{eq:y1_p4}
\end{equation}
%
Inserindo \(y=1\) em \eqref{eq:fp1},
%
\begin{equation}
    x^2+1-1-z=0,
\end{equation}
%
de onde
%
\begin{equation}
    z=x^2.
    \label{eq:z_x2_p4}
\end{equation}
%
Portanto, existe a linha de pontos fixos
%
\begin{equation}
    \boxed{
    \mathcal F:
    \qquad
    (x,y,z)=(x,1,x^2),
    \qquad
    x\ge0.
    }
    \label{eq:fixed_line}
\end{equation}
%
Para o subcaso \(x=0\), a condição \(x'=x\) é automaticamente satisfeita, pois
\(x'=(1+y+z)\cdot0/\mathcal D=0\). A condição \(y'=y\) dá
%
\begin{equation}
    y
    =
    \frac{2y}{1+y^2}.
\end{equation}
%
Multiplicando por \(1+y^2\),
%
\begin{equation}
    y(1+y^2)=2y.
\end{equation}
%
Se \(y=0\), a equação é satisfeita. Se \(y\neq0\), podemos dividir por
\(y\):
%
\begin{equation}
    1+y^2=2
    \quad\Longrightarrow\quad
    y^2=1
    \quad\Longrightarrow\quad
    y=\pm1.
\end{equation}
%
No domínio \(y\ge0\), obtemos
%
\begin{equation}
    y=0
    \qquad\text{ou}\qquad
    y=1.
\end{equation}
%
Para o subcaso \(x=0\), \(y=0\), a condição \(z'=z\) fornece
%
\begin{equation}
    z
    =
    \frac{z^2}{1}
    =
    z^2.
\end{equation}
%
Logo,
%
\begin{equation}
    z(z-1)=0,
\end{equation}
%
de modo que
%
\begin{equation}
    z=0
    \qquad\text{ou}\qquad
    z=1.
\end{equation}
%
Obtemos assim os pontos fixos,
%
\begin{equation}
    \boxed{
    (0,0,0),
    \qquad
    (0,0,1).
    }
\end{equation}
%
Para o subcaso \(x=0\), \(y=1\), a condição \(z'=z\) dá
%
\begin{equation}
    z
    =
    \frac{z^2}{1+1}
    =
    \frac{z^2}{2}.
\end{equation}
%
Multiplicando por \(2\),
%
\begin{equation}
    2z=z^2
    \quad\Longrightarrow\quad
    z(z-2)=0.
\end{equation}
%
Logo,
%
\begin{equation}
    z=0
    \qquad\text{ou}\qquad
    z=2.
\end{equation}
%
O primeiro ponto,
%
\begin{equation}
    \boxed{(0,1,0)},
\end{equation}
%
já pertence à linha \eqref{eq:fixed_line} (basta tomar \(x=0\)).
O segundo,
%
\begin{equation}
    \boxed{
    (0,1,2),
    }
\end{equation}
%
é um ponto fixo adicional da transformação matemática. Entretanto, como as variáveis definidas em
\eqref{eq:x_spin1}--\eqref{eq:z_spin1} satisfazem \(0<x,y,z\le1\)
para \(J,K,D\ge0\), o ponto \((0,1,2)\) não pertence ao domínio físico
correspondente a esses parâmetros. Assim, os principais pontos fixos
matemáticos são
%
\begin{equation}
    \boxed{
    (0,0,0),\qquad
    (0,0,1),\qquad
    (0,1,2),
    \qquad
    (x,1,x^2).
    }
    \label{eq:all_fixed_p4}
\end{equation}
% 
O ponto
%
\begin{equation}
    (1,1,1)
\end{equation}
%
pertence à linha \eqref{eq:fixed_line}. Como
%
\begin{equation}
    x\to1,\qquad
    y\to1,\qquad
    z\to1
\end{equation}
%
corresponde ao limite \(\beta\to0\), ele representa o ponto de alta temperatura, no qual as interações efetivas desaparecem. Já o ponto
%
\begin{equation}
    \boxed{
    (0,0,0)
    }
\end{equation}
%
corresponde ao limite de baixa temperatura no qual os pesos dos
estados excitados são fortemente suprimidos, sendo associado ao regime ordenado. O ponto \((0,0,1)\) pertence ao espaço ampliado de pesos e representa
um regime de competição diferente entre os setores de estados. Ele não é obtido genericamente da subvariedade física inicial \(z=x^2\), mas surge naturalmente quando a transformação é estudada no espaço de Hamiltonianas completo. O ponto \((0,1,2)\) é um ponto fixo da transformação algébrica, mas, como \(z>1\), não pode ser associado diretamente a parâmetros
\(J,K,D\ge0\) através das definições do problema. Sua existência mostra, contudo, que o espaço matemático da transformação contém setores que
não correspondem diretamente ao domínio ferromagnético físico
considerado inicialmente. A linha de pontos fixos
%
\begin{equation}
    y=1,\qquad z=x^2
\end{equation}
%
corresponde a \(J=0\), pois
%
\begin{equation}
    y=e^{-2\beta J}=1
    \quad\Longrightarrow\quad
    J=0.
\end{equation}
% 
A estabilidade pode ser estudada pela matriz jacobiana
%
\begin{equation}
    \mathcal J
    =
    \frac{\partial(x',y',z')}
         {\partial(x,y,z)}.
    \label{eq:jacobian_p4}
\end{equation}
%
Calculamos as derivadas parciais das funções de recorrência
\eqref{eq:xrec}--\eqref{eq:zrec}. Seja
%
\begin{equation}
    \mathcal D=1+x^2+y^2.
\end{equation}
%
Então
%
\begin{align}
    \pd{x'}{x}
    &=
    \frac{(1+y+z)(\mathcal D-2x^2)}{\mathcal D^2},
    \label{eq:dxdx_p4}
    \\[4pt]
    \pd{x'}{y}
    &=
    \frac{x[\mathcal D-2y(1+y+z)]}{\mathcal D^2},
    \label{eq:dxdy_p4}
    \\[4pt]
    \pd{x'}{z}
    &=
    \frac{x}{\mathcal D},
    \label{eq:dxdz_p4}
    \\[4pt]
    \pd{y'}{x}
    &=
    \frac{2x[\mathcal D-(2y+x^2)]}{\mathcal D^2},
    \label{eq:dydx_p4}
    \\[4pt]
    \pd{y'}{y}
    &=
    \frac{2\mathcal D-2y(2y+x^2)}{\mathcal D^2},
    \label{eq:dydy_p4}
    \\[4pt]
    \pd{y'}{z}
    &=
    0,
    \label{eq:dydz_p4}
    \\[4pt]
    \pd{z'}{x}
    &=
    \frac{2x[2\mathcal D-(z^2+2x^2)]}{\mathcal D^2},
    \label{eq:dzdx_p4}
    \\[4pt]
    \pd{z'}{y}
    &=
    -\frac{2y(z^2+2x^2)}{\mathcal D^2},
    \label{eq:dzdy_p4}
    \\[4pt]
    \pd{z'}{z}
    &=
    \frac{2z}{\mathcal D}.
    \label{eq:dzdz_p4}
\end{align}
%
Avaliamos \(\mathcal J\) em cada ponto fixo. Para o ponto \((1,1,1)\), \(\mathcal D=1+1+1=3\). Substituindo nas expressões acima,
%
\begin{align}
    \pd{x'}{x}
    &=
    \frac{(1+1+1)(3-2\cdot1)}{9}
    =
    \frac{3\cdot1}{9}
    =
    \frac13,
    \\[4pt]
    \pd{x'}{y}
    &=
    \frac{1[3-2\cdot1\cdot3]}{9}
    =
    \frac{3-6}{9}
    =
    -\frac13,
    \\[4pt]
    \pd{x'}{z}
    &=
    \frac{1}{3},
    \\[4pt]
    \pd{y'}{x}
    &=
    \frac{2\cdot1[3-(2+1)]}{9}
    =
    0,
    \\[4pt]
    \pd{y'}{y}
    &=
    \frac{2\cdot3-2\cdot1(2+1)}{9}
    =
    0,
    \\[4pt]
    \pd{y'}{z}
    &=
    0,
    \\[4pt]
    \pd{z'}{x}
    &=
    \frac{2\cdot1[2\cdot3-(1+2)]}{9}
    =
    \frac{2\cdot3}{9}
    =
    \frac23,
    \\[4pt]
    \pd{z'}{y}
    &=
    -\frac{2\cdot1(1+2)}{9}
    =
    -\frac23,
    \\[4pt]
    \pd{z'}{z}
    &=
    \frac{2\cdot1}{3}
    =
    \frac23.
\end{align}
%
Logo,
%
\begin{equation}
    \mathcal J(1,1,1)
    =
    \begin{pmatrix}
        \frac13 & -\frac13 & \frac13 \\[4pt]
        0 & 0 & 0 \\[4pt]
        \frac23 & -\frac23 & \frac23
    \end{pmatrix}.
\end{equation}
%
Os autovalores são obtidos de
%
\begin{equation}
    \det(\mathcal J-\lambda I)=0.
\end{equation}
%
Expandindo,
%
\begin{equation}
    \det
    \begin{pmatrix}
        \frac13-\lambda & -\frac13 & \frac13 \\[4pt]
        0 & -\lambda & 0 \\[4pt]
        \frac23 & -\frac23 & \frac23-\lambda
    \end{pmatrix}
    =
    (-\lambda)
    \left[
        \left(\frac13-\lambda\right)
        \left(\frac23-\lambda\right)
        -
        \frac29
    \right]
    =
    0.
\end{equation}
%
O termo entre colchetes é
%
\begin{equation}
    \frac{2}{9}
    -
    \lambda
    +
    \lambda^2
    -
    \frac{2}{9}
    =
    \lambda^2-\lambda
    =
    \lambda(\lambda-1).
\end{equation}
%
Portanto,
%
\begin{equation}
    -\lambda\cdot\lambda(\lambda-1)=0
    \quad\Longrightarrow\quad
    \lambda=0,0,1.
\end{equation}
%
Assim,
%
\begin{equation}
    \lambda=(1,0,0).
\end{equation}
%
O autovalor unitário está associado à direção tangente da linha de pontos fixos, enquanto as outras duas direções são contraídas. Para o ponto \((0,0,0)\), \(\mathcal D=1\). Substituindo,
%
\begin{align}
    \pd{x'}{x} &= 1, &
    \pd{x'}{y} &= 0, &
    \pd{x'}{z} &= 0,
    \\
    \pd{y'}{x} &= 0, &
    \pd{y'}{y} &= 2, &
    \pd{y'}{z} &= 0,
    \\
    \pd{z'}{x} &= 0, &
    \pd{z'}{y} &= 0, &
    \pd{z'}{z} &= 0.
\end{align}
%
Logo,
%
\begin{equation}
    \mathcal J(0,0,0)
    =
    \begin{pmatrix}
        1 & 0 & 0 \\
        0 & 2 & 0 \\
        0 & 0 & 0
    \end{pmatrix},
\end{equation}
%
cujos autovalores são
%
\begin{equation}
    \lambda=(2,1,0).
\end{equation}
%
Para o ponto \((0,0,1)\), \(\mathcal D=1\). Substituindo,
%
\begin{align}
    \pd{x'}{x} &= (1+0+1)(1)=2, &
    \pd{x'}{y} &= 0, &
    \pd{x'}{z} &= 0,
    \\
    \pd{y'}{x} &= 0, &
    \pd{y'}{y} &= 2, &
    \pd{y'}{z} &= 0,
    \\
    \pd{z'}{x} &= 0, &
    \pd{z'}{y} &= 0, &
    \pd{z'}{z} &= 2.
\end{align}
%
Portanto,
%
\begin{equation}
    \mathcal J(0,0,1)
    =
    \begin{pmatrix}
        2 & 0 & 0 \\
        0 & 2 & 0 \\
        0 & 0 & 2
    \end{pmatrix},
\end{equation}
%
e os autovalores são
%
\begin{equation}
    \lambda=(2,2,2).
\end{equation}
%
Para o ponto \((0,1,2)\), \(\mathcal D=1+0+1=2\). Substituindo,
%
\begin{align}
    \pd{x'}{x} &= (1+1+2)(2)/4 = 2, &
    \pd{x'}{y} &= 0, &
    \pd{x'}{z} &= 0,
    \\
    \pd{y'}{x} &= 0, &
    \pd{y'}{y} &= [2\cdot2-2\cdot1(2+0)]/4 = 0, &
    \pd{y'}{z} &= 0,
    \\
    \pd{z'}{x} &= 0, &
    \pd{z'}{y} &= -2\cdot1(4+0)/4 = -2, &
    \pd{z'}{z} &= 2\cdot2/2 = 2.
\end{align}
%
Logo,
%
\begin{equation}
    \mathcal J(0,1,2)
    =
    \begin{pmatrix}
        2 & 0 & 0 \\
        0 & 0 & 0 \\
        0 & -2 & 2
    \end{pmatrix}.
\end{equation}
%
O polinômio característico é
%
\begin{equation}
    \det(\mathcal J-\lambda I)
    =
    (2-\lambda)
    \det
    \begin{pmatrix}
        -\lambda & 0 \\
        -2 & 2-\lambda
    \end{pmatrix}
    =
    (2-\lambda)(-\lambda)(2-\lambda)
    =
    -\lambda(2-\lambda)^2.
\end{equation}
%
Portanto,
%
\begin{equation}
    \boxed{
    \lambda=(2,2,0).
    }
\end{equation}
%
Esses autovalores caracterizam as direções relevantes e irrelevantes do fluxo de renormalização em torno de cada ponto.
%
\end{document}
